% Part: first-order-logic
% Chapter: syntax-and-semantics
% Section: formation-sequences

\documentclass[../../../include/open-logic-section]{subfiles}

\begin{document}

\olfileid{fol}{syn}{fseq}

\olsection{रचनाक्रमिका}

!!{formula}s ची विगमनाधारित व्याख्या देणे आणि त्याला पूरक तंत्र म्हणून
!!{formula}s चे गुणधर्म विगमनाने सिद्ध करणे हा सुबक व कार्यक्षम
दृष्टिकोन आहे. तथापि, !!{formula}s च्या रचनेचा तळापासून वर जाणारा,
टप्प्याटप्प्याचा दृष्टिकोनही उपयुक्त ठरू शकतो; येथे आपण तो
\emph{रचनाक्रमिका} या संकल्पनेच्या साहाय्याने मांडतो.
%
पदे आणि !!{formula}s त्यांच्या विगमनाधारित व्याख्यांचा निर्देश न करता
या रीतीने कशी मांडता येतात हे दाखवण्यासाठी, आपण प्रथम एखाद्या
भाषा~$\Lang L$ मधील चिन्हांपासून बनलेल्या स्वेच्छ चिन्हमालेची संकल्पना
मांडतो.

\begin{defn}[चिन्हमाला]
\ollabel{defn:string}
$\Lang L$ ही प्रथम-क्रम भाषा आहे असे समजा. \emph{$\Lang L$-चिन्हमाला}
म्हणजे $\Lang L$ च्या चिन्हांची सांत क्रमिका. संदर्भावरून भाषा~$\Lang L$
स्पष्टपणे ठरलेली असेल, तर $\Lang L$-चिन्हमालेला आपण अनेकदा फक्त
\emph{चिन्हमाला} म्हणू.
\end{defn}

\begin{ex}
कोणत्याही प्रथम-क्रम भाषा $\Lang L$ साठी सर्व $\Lang L$-!!{formula}s
या $\Lang L$-चिन्हमाला असतात; पण याचे उलट खरे नसते. उदाहरणार्थ,
\[)(\Obj v_0\lif\lexists\] ही $\Lang L$-चिन्हमाला आहे, पण
$\Lang L$-!!{formula} नाही.
\end{ex}

\begin{defn}[पदांसाठी रचनाक्रमिका]
\ollabel{defn:fseq-trm}
$\Lang L$-चिन्हमालांची सांत क्रमिका $\tuple{t_0,\dotsc,t_n}$ ही पद~$t$
साठी \emph{रचनाक्रमिका} असते, जर $t \ident t_n$ आणि प्रत्येक $i \leq n$
साठी एकतर $t_i$ हे !!a{variable} किंवा !!a{constant} असेल, अथवा $\Lang
L$ मध्ये एखादे $k$-स्थानी !!{function}~$f$ असेल आणि असे
$m_1,\dotsc,m_k < i$ अस्तित्वात असतील की
$t_i \ident f(t_{m_1},\dotsc,t_{m_k})$. फरक स्पष्ट करणे आवश्यक असेल,
तर पदांसाठीच्या रचनाक्रमिकांना आपण \emph{पद-रचनाक्रमिका} म्हणू.
%READERNOTE{OLFOL-006: गोठवलेल्या इंग्रजी व्याख्येत k-स्थानी फलनासाठी m_0 ते m_k असे k+1 निर्देशांक आणि तितकेच युक्तिवाद दिले आहेत. येथे k युक्तिवादांसाठी m_1 ते m_k ही सुसंगत मांडणी वापरली आहे.}
\end{defn}

\begin{ex}
क्रमिका
\[
    \tuple{\Obj c_0, \Obj v_0, \Atom{\Obj f^2_0}{\Obj c_0, \Obj v_0}, \Atom{\Obj f^1_0}{\Atom{\Obj f^2_0}{\Obj c_0, \Obj v_0}}}
\]
ही पद $\Atom{\Obj f^1_0}{\Atom{\Obj
f^2_0}{\Obj c_0, \Obj v_0}}}$ साठी रचनाक्रमिका आहे; त्याचप्रमाणे
%READERNOTE{OLFOL-009: या पदाच्या गोठवलेल्या इंग्रजी आवृत्तीत बाहेरील फलन-अनुप्रयोगाचा शेवटचा आकुंचित कंस राहिला आहे. येथे तो कंस पूर्ण केला आहे.}
\[
    \tuple{\Obj v_0, \Obj c_0, \Atom{\Obj f^2_0}{\Obj c_0, \Obj v_0}, \Atom{\Obj f^1_0}{\Atom{\Obj f^2_0}{\Obj c_0, \Obj v_0}}}.
\]
हीदेखील त्याची रचनाक्रमिका आहे.
\end{ex}

\begin{defn}[सूत्रांसाठी रचनाक्रमिका]
\ollabel{defn:fseq-frm}
$\Lang L$-चिन्हमालांची सांत क्रमिका $\tuple{!A_0,\dotsc,!A_n}$ ही
$!A$ साठी \emph{रचनाक्रमिका} असते, जर $!A \ident !A_n$ आणि प्रत्येक
$i \leq n$ साठी एकतर $!A_i$ हे आण्विक !!{formula} असेल, अथवा पुढीलपैकी
एखादी अट पूर्ण करणारे $j,k < i$ आणि एखादे !!a{variable}~$x$ अस्तित्वात
असतील:
\begin{enumerate}
    \tagitem{prvNot}{$!A_i \ident \lnot !A_j$.}{}%
    \tagitem{prvAnd}{$!A_i \ident (!A_j \land !A_k)$.}{}%
    \tagitem{prvOr}{$!A_i \ident (!A_j \lor !A_k)$.}{}%
    \tagitem{prvIf}{$!A_i \ident (!A_j \lif !A_k)$.}{}%
    \tagitem{prvIff}{$!A_i \ident (!A_j \liff !A_k)$.}{}%
    \tagitem{prvAll}{$!A_i \ident \lforall[x][!A_j]$.}{}%
    \tagitem{prvEx}{$!A_i \ident \lexists[x][!A_j]$.}{}%
\end{enumerate}
फरक स्पष्ट करणे आवश्यक असेल, तर सूत्रांसाठीच्या रचनाक्रमिकांना आपण
\emph{सूत्र-रचनाक्रमिका} म्हणू.
\end{defn}

\begin{ex}
\[
\tuple{
    \Atom{\Obj A^1_0}{\Obj v_0},
    \Atom{\Obj A^1_1}{\Obj c_1},
    (\Atom{\Obj A^1_1}{\Obj c_1} \land \Atom{\Obj A^1_0}{\Obj v_0}),
    \lexists[\Obj v_0][(\Atom{\Obj A^1_1}{\Obj c_1} \land \Atom{\Obj A^1_0}{\Obj v_0})]
}
\]
ही $\lexists[\Obj v_0][(\Atom{\Obj A^1_1}{\Obj c_1} \land \Atom{\Obj
A^1_0}{\Obj v_0})]$ ची एक रचनाक्रमिका आहे; त्याचप्रमाणे
\begin{multline*}
\tuple{
    \Atom{\Obj A^1_0}{\Obj v_0},
    \Atom{\Obj A^1_1}{\Obj c_1},
    (\Atom{\Obj A^1_1}{\Obj c_1} \land \Atom{\Obj A^1_0}{\Obj v_0}),
    \Atom{\Obj A^1_1}{\Obj c_1},\\
    \lforall[\Obj v_1][\Atom{\Obj A^1_0}{\Obj v_0}],
    \lexists[\Obj v_0][(\Atom{\Obj A^1_1}{\Obj c_1} \land \Atom{\Obj A^1_0}{\Obj v_0})]
}.
\end{multline*}
%
दुसऱ्या उदाहरणावरून दिसते की रचनाक्रमिकेत ``अडगळ'' असू शकते:
अनावश्यक असलेली किंवा रचनेत काहीही हातभार न लावणारी !!{formula}s.
\end{ex}

\begin{prop}\ollabel{prop:formed}
$\Frm[L]$ मधील प्रत्येक !!{formula}~$!A$ ला एक रचनाक्रमिका असते.
\end{prop}

\begin{proof}
$!A$ आण्विक आहे असे समजा. मग क्रमिका~$\tuple{!A}$ ही $!A$ साठी
रचनाक्रमिका आहे.
%
आता $!B$ आणि~$!C$ यांना अनुक्रमे $\tuple{!B_0,\dotsc,!B_n}$ आणि
$\tuple{!C_0,\dotsc,!C_m}$ या रचनाक्रमिका आहेत असे समजा.
%
\begin{enumerate}
  \tagitem{prvNot}{जर $!A \ident \lnot !B$, तर
      $\tuple{!B_0,\dotsc,!B_n,\lnot !B_n}$ ही $!A$ साठी
      रचनाक्रमिका आहे.}{}
  \tagitem{prvAnd}{जर $!A \ident (!B \land !C)$, तर
      $\tuple{!B_0,\dotsc,!B_n,!C_0,\dotsc,!C_m,(!B_n \land !C_m)}$
      ही $!A$ साठी रचनाक्रमिका आहे.}{}
  \tagitem{prvOr}{जर $!A \ident (!B \lor !C)$, तर
      $\tuple{!B_0,\dotsc,!B_n,!C_0,\dotsc,!C_m,(!B_n \lor !C_m)}$
      ही $!A$ साठी रचनाक्रमिका आहे.}{}
  \tagitem{prvIf}{जर $!A \ident (!B \lif !C)$, तर
      $\tuple{!B_0,\dotsc,!B_n,!C_0,\dotsc,!C_m,(!B_n \lif !C_m)}$
      ही $!A$ साठी रचनाक्रमिका आहे.}{}
  \tagitem{prvIff}{जर $!A \ident (!B \liff !C)$, तर
      $\tuple{!B_0,\dotsc,!B_n,!C_0,\dotsc,!C_m,(!B_n \liff !C_m)}$
      ही $!A$ साठी रचनाक्रमिका आहे.}{}
  \tagitem{prvAll}{जर $!A \ident \lforall[x][!B]$, तर
      $\tuple{!B_0,\dotsc,!B_n,\lforall[x][!B_n]}$
      ही $!A$ साठी रचनाक्रमिका आहे.}{}
  \tagitem{prvEx}{जर $!A \ident \lexists[x][!B]$, तर
      $\tuple{!B_0,\dotsc,!B_n,\lexists[x][!B_n]}$
      ही $!A$ साठी रचनाक्रमिका आहे.}{}
\end{enumerate}
!!{formula}s वरील विगमनाच्या तत्त्वानुसार प्रत्येक !!{formula} ला एक
रचनाक्रमिका असते.
\end{proof}

याचे उलट विधानही आपण सिद्ध करू शकतो. ते महत्त्वाचे आहे, कारण त्यामुळे
सूत्रे परिभाषित करण्याच्या आपल्या दोन्ही पद्धती सममूल्य आहेत, म्हणजे
त्या समान परिणाम देतात, हे दिसते. तसेच रचनाक्रमिकांच्या लांबीवर साधे
विगमन वापरून सूत्रांविषयीची प्रमेये सिद्ध करता येतात.

\begin{lem}
\ollabel{lem:fseq-init}
$\tuple{!A_0,\dotsc,!A_n}$ ही $!A_n$ साठी रचनाक्रमिका आहे आणि
$k \leq n$ असे समजा. मग $\tuple{!A_0,\dotsc,!A_k}$ ही $!A_k$ साठी
रचनाक्रमिका आहे.
\end{lem}

\begin{proof}
अभ्यासप्रश्न.
\end{proof}

\begin{prob}
\olref[fol][syn][fseq]{lem:fseq-init} सिद्ध करा.
\end{prob}

\begin{thm}
\ollabel{thm:fseq-frm-equiv}
$\Frm[L]$ हा अशा सर्व $\Lang L$-चिन्हमाला~$!A$ चा संच आहे की $!A$
साठी एखादी सूत्र-रचनाक्रमिका अस्तित्वात आहे.
\end{thm}

\begin{proof}
भाषा~$\Lang L$ मधील ज्या सर्व चिन्हमालांना रचनाक्रमिका आहे, त्यांचा
संच~$F$ असू द्या. $\Frm[L] \subseteq F$ हे
\olref[fol][syn][fseq]{prop:formed} मध्ये आपण पाहिले आहे; म्हणून आता उलट
समावेश सिद्ध करू.

$!A$ ला $\tuple{!A_0,\dotsc,!A_n}$ ही रचनाक्रमिका आहे असे समजा.
$n$ वरील प्रबल विगमनाने $!A \in \Frm[L]$ हे सिद्ध करू. लांबी
$m < n$ असलेली रचनाक्रमिका ज्या प्रत्येक चिन्हमालेला आहे, ती
$\Frm[L]$ मध्ये आहे, हे आपले विगमन गृहीतक आहे. रचनाक्रमिकेच्या
व्याख्येनुसार एकतर $!A \ident !A_n$ आण्विक आहे किंवा पुढीलपैकी एखादी
अट पूर्ण करणारे $j,k < n$ अस्तित्वात असले पाहिजेत:
\begin{enumerate}
    \tagitem{prvNot}{$!A \ident \lnot !A_j$.}{}%
    \tagitem{prvAnd}{$!A \ident (!A_j \land !A_k)$.}{}%
    \tagitem{prvOr}{$!A \ident (!A_j \lor !A_k)$.}{}%
    \tagitem{prvIf}{$!A \ident (!A_j \lif !A_k)$.}{}%
    \tagitem{prvIff}{$!A \ident (!A_j \liff !A_k)$.}{}%
    \tagitem{prvAll}{$!A \ident \lforall[x][!A_j]$.}{}%
    \tagitem{prvEx}{$!A \ident \lexists[x][!A_j]$.}{}%
\end{enumerate}
आता प्रत्येक शक्यतेचा स्वतंत्र विचार करू. $!A$ आण्विक असेल, तर
$!A_n \in \Frm[L]$. त्याऐवजी $!A \ident (!A_j \land !A_k)$ असे समजा.
%READERNOTE{OLFOL-007: हा पुरावा सामान्य भाषा L साठी आहे; गोठवलेल्या इंग्रजीतील या परिच्छेदाच्या दोन ठिकाणी चुकून विधानीय भाषा L_0 आली आहे. येथे प्रमेयाशी सुसंगत L वापरले आहे.}
%READERNOTE{OLFOL-008: गोठवलेल्या इंग्रजीत या एकाच रचना-बाबीत विन्यासात्मक अनन्यता-चिन्हाऐवजी पर्यायी सममूल्यता-चिन्ह आले आहे. व्याख्या आणि इतर सर्व बाबींशी सुसंगत विन्यासात्मक अनन्यता येथे वापरली आहे.}
\olref[fol][syn][fseq]{lem:fseq-init} नुसार
$\tuple{!A_0,\dotsc,!A_j}$ आणि $\tuple{!A_0,\dotsc,!A_k}$ या अनुक्रमे
$!A_j$ आणि~$!A_k$ साठी रचनाक्रमिका आहेत. या $!A$ च्या
रचनाक्रमिकेच्या उचित आरंभीच्या उपक्रमिका असल्यामुळे दोन्हींची लांबी
$n$ पेक्षा कमी आहे. म्हणून विगमन गृहीतकानुसार $!A_j$ आणि~$!A_k$ ही
$\Frm[L]$ मध्ये आहेत; त्यामुळे !!a{formula} च्या व्याख्येनुसार
$(!A_j \land !A_k)$ हेही त्यात आहे. इतर शक्यता समांतर युक्तिवादाने
सिद्ध होतात.
\end{proof}

पदांसाठीच्या रचनाक्रमिकांना !!{formula}s च्या रचनाक्रमिकांसारखेच गुणधर्म
असतात.

\begin{prop}
\ollabel{prop:fseq-trm-equiv}
$\Trm[L]$ हा अशा सर्व $\Lang L$-चिन्हमाला $t$ चा संच आहे की $t$
साठी एखादी पद-रचनाक्रमिका अस्तित्वात आहे.
\end{prop}

\begin{proof}
अभ्यासप्रश्न.
\end{proof}

\begin{prob}
\olref[fol][syn][fseq]{prop:fseq-trm-equiv} सिद्ध करा.
सूचना: \olref[fol][syn][fseq]{thm:fseq-frm-equiv} च्या सिद्धतेत वापरलेली
रणनीती वापरा.
\end{prob}

रचनाक्रमिकेत दोन प्रकारची ``अडगळ'' येऊ शकते: पुनरावृत्त घटक आणि सूत्र
किंवा पद यांच्या रचनेशी असंबद्ध घटक. न्यूनतम रचनाक्रमिका विचारात घेऊन
आपण दोन्ही प्रकार दूर करू शकतो.

\begin{defn}[न्यूनतम रचनाक्रमिका]
\ollabel{defn:minimal-fseq}
सूत्र~$!A$ साठीची रचनाक्रमिका $\tuple{!A_0, \ldots, !A_n}$ ही
$!A$ साठी \emph{न्यूनतम रचनाक्रमिका} असते, जर $!A$ साठीची प्रत्येक
दुसरी रचनाक्रमिका~$s$ घेतली असता~$s$ ची लांबी $n+1$ पेक्षा अधिक किंवा
तिच्याइतकी असेल.

त्याचप्रमाणे, पद~$t$ साठीची रचनाक्रमिका $\tuple{t_0, \ldots, t_n}$ ही
$t$ साठी \emph{न्यूनतम रचनाक्रमिका} असते, जर $t$ साठीची प्रत्येक
दुसरी रचनाक्रमिका~$s$ घेतली असता~$s$ ची लांबी $n+1$ पेक्षा अधिक किंवा
तिच्याइतकी असेल.
\end{defn}

एका सूत्राला किंवा पदाला एकापेक्षा अधिक न्यूनतम रचनाक्रमिका असू शकतात;
पण त्यांत नेमक्या त्याच चिन्हमाला असतील, हे लक्षात घ्या.

\begin{prop}
\ollabel{prop:subformula-equivs}
पुढील विधाने सममूल्य आहेत:
\begin{enumerate}
    \item $!B$ हे $!A$ चे उप-!!{formula} आहे.
    \item $!B$ हे $!A$ च्या प्रत्येक रचनाक्रमिकेत येते.
    \item $!B$ हे $!A$ च्या एखाद्या न्यूनतम रचनाक्रमिकेत येते.
\end{enumerate}
\end{prop}

\begin{proof}
अभ्यासप्रश्न.
\end{proof}

\begin{prob}
\olref[fol][syn][fseq]{prop:subformula-equivs} सिद्ध करा.
\end{prob}

\begin{history}
रेमंड स्मलियन यांनी त्यांच्या \emph{First-Order Logic}
\citep{Smullyan1968} या पाठ्यपुस्तकात रचनाक्रमिका प्रथम मांडल्या.
रचनाक्रमिकांचे आणखी गुणधर्म \citet{Zuckerman1973} यांनी प्रस्थापित केले.
\end{history}

\end{document}
